\section*{Chapter \ref{ch:b1:national-accounts-in-depth} --- National Accounts in Depth}
% Numbers recomputed in code/bachelor-1/18-national-accounts-in-depth/exercises/solutions.py (tested).

\begin{solution}{exo:b1:national-accounts-in-depth:gni-ex}
$500 - 20 = 480$.
\end{solution}

\begin{solution}{exo:b1:national-accounts-in-depth:three}
$30 - 10 - 35 = -15$: a current-account deficit of 15, and the rest of the world
lends 15.
\end{solution}

\begin{solution}{exo:b1:national-accounts-in-depth:read}
Growth rates are in log points a year over 29 years: Mali
$e^{0.01495 \times 29} = 1.54$ against $e^{0.0617 \times 29} = 5.99$; Honduras $5.34$ against $1.43$. The two
sources disagree on whether income per person rose by half or sixfold.
\end{solution}

\begin{solution}{exo:b1:national-accounts-in-depth:recession}
Tax receipts fall and benefits rise, so the government's \omterm{def:b1:national-accounts-in-depth:net-lending}{net lending} falls
(its deficit grows). Households and firms cut investment more than saving,
so the private balance rises; part of the change may also show up in the
current account.
\end{solution}

\begin{solution}{exo:b1:national-accounts-in-depth:chain-ex}
At year-0 prices: 200, 310, 520. Chained: $200 \times 310/200 = 310$ in year 1; from year 1
to 2 at year-1 prices, $(120 \times 1.1 + 40 \times 6)/(110 \times 1.1 + 20 \times 6) = 372/241$, so $310 \times 372/241 = 478.51$.
The chained components are A: $100 \times 1.2 = 120$ and B: $100 \times 4 = 400$ in year-0 prices,
summing to 520, not 478.51.
\end{solution}

\begin{solution}{exo:b1:national-accounts-in-depth:ppp-ex}
The basket costs $200 + 200 + 300 = 700$ dalens and $500 + 300 + 900 = 1700$ kests: a PPP of
$1700/700 = 2.4286$ kests per dalen. Kestria's GDP is $60000/2.4286 = 24{,}706$ million
dalens at PPP and $60000/3 = 20{,}000$ million at the market rate. Its price level
relative to Brenland is $2.4286/3 = 0.81$: things cost $19\%$ less there.
\end{solution}

\begin{solution}{exo:b1:national-accounts-in-depth:brenland}
\omterm{prop:b1:models-and-data:real-growth}{Real growth} $292390/286895 - 1 = 1.92\%$; deflator inflation $134.5/132.3 - 1 = 1.66\%$; nominal
growth $393264/379562 - 1 = 3.61\%$; and $1.0192 \times 1.0166 = 1.0361$.
\end{solution}

\begin{solution}{exo:b1:national-accounts-in-depth:ip}
GDP rises: the subsidiary's value added now includes the returns to the
patents. \omterm{def:b1:national-accounts-in-depth:gni}{GNI} rises much less, since the profits belong to the foreign
parent and are paid out as primary income to the rest of the world.
Employment barely changes.
\end{solution}

\begin{solution}{exo:b1:national-accounts-in-depth:pli-ex}
$20/50 = 0.4$: a basket costing one dollar in the base economy costs the
equivalent of 40 cents there at the market rate.
\end{solution}

\begin{solution}{exo:b1:national-accounts-in-depth:sum}
For each domestic sector, \omterm{def:b1:national-accounts-in-depth:net-lending}{net lending} is saving minus investment plus net
capital transfers received. Summing over sectors, transfers between
domestic sectors cancel, leaving national saving minus investment plus net
capital transfers from abroad. National saving is $Y + \mathrm{NPI} + \mathrm{TR} - C - G$, so the
sum is $X - M + \mathrm{NPI} + \mathrm{TR}$ plus net capital transfers from abroad: the current
account plus the capital account, which is minus the rest of the world's
\omterm{def:b1:national-accounts-in-depth:net-lending}{net lending}.
\end{solution}

\begin{solution}{exo:b1:national-accounts-in-depth:fixed-base}
Let good 1's relative price fall, $p_1^1/p_2^1 < p_1^0/p_2^0$, and its quantity grow faster,
$q_1^1/q_1^0 > q_2^1/q_2^0$. Growth at base prices is a weighted mean of $q_1^1/q_1^0$ and
$q_2^1/q_2^0$ with weights $p_i^0q_i^0$; at current prices, with weights $p_i^1q_i^0$. The second
puts relatively less weight on good 1, the fast-growing one, so it is
smaller. (A chained index updates the weights every period; it lies
between the two when relative prices move the same way in every period,
the remark after the proposition.)
\end{solution}

\begin{solution}{exo:b1:national-accounts-in-depth:refute}
Quantities that grow faster than measured output: electricity use, cargo,
night lights seen from space, phone subscriptions, tax receipts; household
surveys of consumption growing faster than the accounts' consumption; a
base year decades old; a large informal sector that the sources miss.
Slow growth with all of these slow too is slow growth.
\end{solution}

\begin{solution}{pb:b1:national-accounts-in-depth:rebase}
\textbf{1.} $225014 + 88484 + 78653 + 129777 - 128664 = 393264$.

\textbf{2.} $393264 - 5900 = 387364$.

\textbf{3.} $387364 - 59000 = 328364$.

\textbf{4.} $387364 + 1200 = 388564$.

\textbf{5.} $388564 - 225014 - 78653 = 84897$.

\textbf{6.} $84897/388564 = 21.85\%$.

\textbf{7.} $157306 - 165171 = -7865$, $-2.0\%$ of GDP.

\textbf{8.} $129777 - 128664 - 5900 + 1200 = -3587$, $-0.91\%$ of GDP.

\textbf{9.} $84897 - 88484 = -3587$.

\textbf{10.} $-3587 - (-7865) = 4278$.

\textbf{11.} The rest of the world lends 3,587: $4278 - 7865 + 3587 = 0$.

\textbf{12.} $236818/393264 = 60.2\%$.

\textbf{13.} $393264/248438 = 1.5829$.

\textbf{14.} $292390/248438 = 1.1769$, or $1.1769^{1/10} - 1 = 1.64\%$ a year.

\textbf{15.} $100 \times 393264/292390 = 134.5$.

\textbf{16.} 1000, 1350 and 2600.

\textbf{17.} 1000, 1350 and $1350 \times 1740/1205 = 1949.4$.

\textbf{18.} $1600 + 1000 = 2600 \ne 1949.4$.

\textbf{19.} The chained volume: it weights each period's growth by
recent prices, so computers, whose price fell by a factor of more than six,
do not keep their 2015 weight. The fixed-base series says the economy grew
$160\%$ when, weighted by what goods cost when they were produced, it grew
about $95\%$.

\textbf{20.} $248438 \times 1.128 = 280238.1$ and $292390 \times 1.14 = 333324.6$.

\textbf{21.} By $14\%$.

\textbf{22.} $(333324.6/280238.1)^{1/10} - 1 = 1.75\%$ a year, against $1.64\%$: $0.11$ points more.

\textbf{23.} $236818/(393264 \times 1.14) = 52.8\%$.

\textbf{24.} The informal sector was there in both years: adding it raises
both levels by similar shares, and growth depends only on the ratio of the
two, which moves from 1.128 to 1.14.

\textbf{25.} The rebasing raises 2025 GDP by $14\%$ but average growth from 2015 to
2025 by only $0.11$ points a year.
\end{solution}
